Castalia is a four-hart (four-core) multiprocessor built from four identical VestaRV cores. Each hart is a complete, single-issue, in-order RV32IMAC processor with its own private memory, so the common case --- instruction fetch and private data access --- never contends with the other harts and runs at full core clock. The harts cooperate through a single arbitrated \textit{shared window} in the address space, which contains a small shared RAM, the inter-processor interrupt hardware (\peripheral{CLINT}), a hardware mutex bank, an interrupt router, and shared-access aliases of most peripherals. The main features of the multi-core architecture are listed below:

\begin{itemize}
	\item Four identical VestaRV harts, each identified by the read-only \register{mhartid} CSR (address \texttt{0xF14})
	\item Private memory per hart: each hart has its own RAM (and hart 0 the boot ROM), so fetch and private loads/stores never stall
	\item One shared window at \texttt{0x10000}--\texttt{0x13FFF}, reachable by all four harts through a serializing round-robin arbiter
	\item 1 KiB of shared RAM for locks, boot mailboxes, and inter-hart data
	\item \peripheral{CLINT} with per-hart software interrupts (inter-processor interrupts) and a shared 64-bit \register{mtime} with per-hart compare registers
	\item Sixteen single-instruction hardware mutexes (\peripheral{MUTEX} bank)
	\item Per-hart peripheral interrupt routing (\peripheral{IRQROUTER}), so any peripheral interrupt can be steered to any hart
	\item \texttt{LR}/\texttt{SC} and \texttt{AMO} atomic instructions supported to shared RAM, including sub-word shared stores (byte lanes)
	\item Parked harts sleep in \asminline{wfi} and are woken by a software interrupt --- no busy-wait power burn
\end{itemize}



\subsection{Harts and Roles}

Hart 0 is the \textit{management hart}. It is the only hart attached to the boot path (SPI flash via \peripheral{SPI0} and the boot ROM) and the only hart wired to the private-only peripherals (\peripheral{GPIO0}, \peripheral{SYSTEM}, \peripheral{SARADC}, \peripheral{AFE}). Its address map is identical to the single-core VestaRV MCU, plus the shared window.

Harts 1--3 are \textit{tile harts}. Each tile contains an unmodified VestaRV core plus its own private RAM0 (\texttt{0x8000}) and RAM1 (\texttt{0xC000}), replicating the proven single-core memory path. Tile harts do not have the SPI-flash boot path; they come out of reset executing from their private RAM (reset PC \texttt{0x8200}) and are expected to park immediately until hart 0 releases them (see Section \ref{ss:multicore-boot}).

Software discovers which hart it is running on by reading the \register{mhartid} CSR:

\begin{verbatim}
    csrr    t0, mhartid      # t0 = 0, 1, 2, or 3
\end{verbatim}



\subsection{The Shared Window}

All cross-hart communication happens through the shared window at \texttt{0x10000}--\texttt{0x13FFF}. Every hart sees the same physical devices at the same addresses:

\begin{center}
\begin{tabular}{ l l l }
	\textbf{Address range} & \textbf{Device} & \textbf{Notes} \\ \hline
	\texttt{0x10000}--\texttt{0x103FF} & Shared RAM (1 KiB) & Locks, mailboxes, inter-hart data \\
	\texttt{0x11000}--\texttt{0x11FFF} & \peripheral{CLINT} & IPIs (\register{MSIPx}), \register{MTIME}/\register{MTIMECMPx} \\
	\texttt{0x12000}--\texttt{0x12FFF} & Shared \peripheral{UART0} & Console UART, shared-access alias \\
	\texttt{0x13000}--\texttt{0x13FFF} & Shared peripheral page & 16 slots of 256 bytes, see below \\ \hline
\end{tabular}
\end{center}

The shared window is for \textbf{data access only}: harts cannot fetch instructions from it. Code must live in each hart's private memory.

\subsubsection{Arbitration and Stalling}

A round-robin arbiter serializes all shared-window transactions. The arbiter grants one hart at a time and holds the grant until that hart's whole transaction completes, then advances the round-robin pointer, so no hart can be starved. While a hart waits for its grant, its clock is stalled transparently --- software never sees a partial transaction and needs no special code for shared accesses. Because a stalled hart contributes nothing until its turn comes, frequently-polled data structures (spin locks in particular) should use a backoff strategy (see Section \ref{ss:multicore-sync}).

An \texttt{AMO} (atomic read-modify-write) instruction holds the arbiter grant across both halves of its read-modify-write pair, so \texttt{AMO}s to shared RAM are atomic with respect to all other harts.

\subsubsection{The Shared Peripheral Page}

The page at \texttt{0x13000}--\texttt{0x13FFF} is divided into sixteen 256-byte slots that deliberately mirror the legacy peripheral page's slot numbering at \texttt{0x4000}: the shared alias of the peripheral in legacy slot $n$ lives at $\mathtt{0x13000} + 256 n$. Through these aliases, harts 1--3 (and hart 0) can operate peripherals with the arbiter guaranteeing register-access atomicity. Two slots are repurposed for multi-core infrastructure blocks: slot 0 holds the \peripheral{MUTEX} bank (its legacy owner \peripheral{GPIO0} is never shared) and slot 9 holds the \peripheral{IRQROUTER} (its legacy owner \peripheral{SYSTEM} remains hart 0's private system controller).

\begin{center}
\begin{tabular}{ l l l l }
	\textbf{Slot} & \textbf{Shared address} & \textbf{Device} & \textbf{Legacy address} \\ \hline
	0  & \texttt{0x13000} & \peripheral{MUTEX} bank (\peripheral{GPIO0} is never shared) & --- \\
	1  & \texttt{0x13100} & \peripheral{GPIO1}  & \texttt{0x4100} \\
	3  & \texttt{0x13300} & \peripheral{SPI1}   & \texttt{0x4300} \\
	5  & \texttt{0x13500} & \peripheral{UART1}  & \texttt{0x4500} \\
	6  & \texttt{0x13600} & \peripheral{TIMER0} & \texttt{0x4600} \\
	7  & \texttt{0x13700} & \peripheral{TIMER1} & \texttt{0x4700} \\
	8  & \texttt{0x13800} & \peripheral{GPIO2}  & \texttt{0x4800} \\
	9  & \texttt{0x13900} & \peripheral{IRQROUTER} (\peripheral{SYSTEM} is never shared) & --- \\
	10 & \texttt{0x13A00} & \peripheral{NPU} (register bus only) & \texttt{0x4A00} \\
	13 & \texttt{0x13D00} & \peripheral{GPIO3}  & \texttt{0x4D00} \\
	14 & \texttt{0x13E00} & \peripheral{I2C0}   & \texttt{0x4E00} \\
	15 & \texttt{0x13F00} & \peripheral{I2C1}   & \texttt{0x4F00} \\ \hline
\end{tabular}
\end{center}

Peripherals not listed (\peripheral{GPIO0}, \peripheral{SYSTEM}, \peripheral{SPI0}, \peripheral{SARADC}, \peripheral{AFE}) are private to hart 0 and are only reachable at their legacy \texttt{0x4000}-page addresses. Note that the \peripheral{NPU}'s shared slot exposes only its register bus; the NPU still operates on vectors in hart 0's private RAM1.



\subsection{Boot and Hart Release} \label{ss:multicore-boot}

Only hart 0 executes the boot ROM and SPI-flash boot sequence, exactly as on the single-core chip. Harts 1--3 come out of reset executing from their private RAM and immediately park, polling their \textit{boot mailbox} in shared RAM with plain loads (no side effects while parked). Hart 0 releases each parked hart by writing its mailbox:

\begin{center}
\begin{tabular}{ l l l }
	\textbf{Mailbox} & \textbf{Address} & \textbf{Meaning} \\ \hline
	\register{ENTRY[h]} & $\mathtt{0x10110} + 4h$ & Entry PC for hart $h$, written \textit{before} \register{GO[h]} \\
	\register{GO[h]}    & $\mathtt{0x10100} + 4h$ & Release: hart 0 writes the magic value \texttt{0x600DB007} \\
	\register{DONE[h]}  & $\mathtt{0x10120} + 4h$ & Hart $h$ writes $\mathtt{0xD00E0000} + h$ after waking \\ \hline
\end{tabular}
\end{center}

The release sequence for hart $h$ is: hart 0 writes the hart's entry point to \register{ENTRY[h]}, then writes the magic value to \register{GO[h]}; hart $h$ observes the magic value, jumps to the entry PC, and acknowledges by writing \register{DONE[h]}. Alternatively (and preferably, once the application is running), a parked hart can sleep in \asminline{wfi} and be woken with a \peripheral{CLINT} software interrupt instead of polling --- see the \peripheral{CLINT} chapter.



\subsection{Synchronization} \label{ss:multicore-sync}

Castalia offers three synchronization mechanisms, from cheapest to most flexible:

\begin{enumerate}
	\item \textbf{Hardware mutexes} (\peripheral{MUTEX} bank): a single load claims a free mutex; a single store releases it. No retry loop hardware state, no reservation. This is the preferred lock for most uses. See the \peripheral{MUTEX} chapter.
	\item \textbf{\texttt{LR}/\texttt{SC} spinlocks in shared RAM}: standard RISC-V reservation-based locks. A failed \texttt{SC} has its write suppressed by the reservation unit, so it is safe under contention.
	\item \textbf{\texttt{AMO} instructions to shared RAM}: atomic add, swap, and friends; the arbiter locks the grant across the read-modify-write pair.
\end{enumerate}

Two rules apply to all of them:

\begin{itemize}
	\item \textbf{Never} use \texttt{LR}/\texttt{SC} or \texttt{AMO} instructions on \peripheral{MUTEX} bank addresses --- hardware mutexes are claimed with plain loads and released with plain stores only.
	\item Under contention, spin with a \textit{hart-scaled backoff}: delay between retries by an amount proportional to \register{mhartid} (for example $\mathtt{delay} = k \cdot (\register{mhartid} + 1)$) and bound the retry count. Four harts hammering the arbiter in lockstep can otherwise livelock a naive retry loop.
\end{itemize}

There are no data caches, so there is no cache-coherence protocol to consider: a shared-RAM write is globally visible as soon as the arbiter completes the transaction.
